\selectlanguage{ngerman}%
\frenchspacing%
\begin{ledgroupsized}[r]{120mm}%
\footnotesize%
\pstart%
\noindent\textbf{Überlieferung:}%
\pend%
\end{ledgroupsized}%
\begin{ledgroupsized}[r]{114mm}%
\footnotesize%
\pstart%
\parindent -6mm%
\makebox[6mm][l]{\textit{L}}%
Aufzeichnung:
LH XXXVII\;5 Bl.\;18.
Ein Blatt, 4\textsuperscript{o};
Ränder beschnitten;
Papiererhaltungsmaßnahmen.
Eine voll beschriebene Seite auf Bl.\;\href{\Scan{LH_37_05_018r}}{18\,r\textsuperscript{o}}.
Bl.\;\href{\Scan{LH_37_05_018v}}{18\,v\textsuperscript{o}} ist unbeschrieben.
\pend%
\end{ledgroupsized}%
\selectlanguage{latin}%
\frenchspacing%
\vspace{8mm}%
\normalsize%
\pstart%
\noindent%
\bracks{\href{\Scan{LH_37_05_018r}}%
{18\,r\textsuperscript{o}}}%
\hfill%
\edlabel{LH_37_05_018_018r_anfg-1}%
Martii 1677%
\edtext{}{%
\xxref{LH_37_05_018_018r_anfg-1}{LH_37_05_018_018r_anfg-2}%
\lemma{1677}%
\Bfootnote{%
\bracks{/}
\textit{(1)}~Ut %
\textit{(2)}~Experimentis cognoscendum est,%
~\textit{L}%
}}%
\pend%
\pstart%
\noindent%
Experimentis%
\indexS{experimentum}
cognoscendum est,%
\edlabel{LH_37_05_018_018r_anfg-2}
quod dudum rationibus%
\indexS{ratio}
conjicio,
\edtext{%
vim in obstaculum%
\indexS{obstaculum}%
\indexS{vis agens in obstaculum}
solum agere,%
}{%
\lemma{vim}%
\Bfootnote{%
\textit{(1)}~non n %
\textit{(2)}~in obstaculum solum agere,%
~\textit{L}%
}}
cum potest non in id
quod ei connexum est,
modo vis connectens%
\indexS{vis connectens tabulas}
vi compellentis cedat;%
\indexS{vis compellentis}
e.\,g.
sint duae
tabulae planae%
\indexS{tabulae planae applicatae sibi}
sibi perfecte applicatae%
\indexS{tabulae planae applicatae sibi}
\textit{A},\;\textit{B} %
horizonti parallelae,
inferior \textit{A} mensae libere imposita,%
\indexS{mensa}
superior ei incumbens \textit{B},
patet
mediocri vi%
\indexS{vis mediocris}%
~%
\edtext{%
\textit{C}%
}{%
\lemma{\textit{C}}%
\Bfootnote{%
\textit{erg.\;L}%
}}
non posse elevari tabulam \textit{B},
quin simul elevetur tabula \textit{A};
si vero \textit{C} sit pondus ingens,%
\indexS{pondus ingens}
quale scilicet aqu\angles{ae}%
\indexS{aqua 30 pedum}
triginta pedum%
\indexS{triginta pes aquae}%
\indexS{pes}
\edtext{%
crassitiei%
\indexS{crassities}
pondus%
}{%
\lemma{crassitiei}%
\Bfootnote{%
\bars{suae \textit{streicht Hrsg.}} pondus %
~\textit{L}%
}}
aequat,
tunc superabit vim tabulas connectentem,%
\indexS{vis connectens tabulas}%
\indexS{tabulae planae applicatae sibi}
adeoque eo casu ajo
solam superiorem tabulam levatum iri
inferiore intacta.
\pend%
\begin{edfigure}[1.750em][1.250em]%
\includegraphics%
{gesamttex/edit_VIII-4/images/LH_37_05_018_d1_018r.pdf}%
\edcaption[0.250em]{%
\bracks{\textit{Fig.\:1}}%
\edtext{}{%
\lemma{\bracks{\textit{Fig.\:1}}\,}%
\Cfootnote[nonum]{%
Zusätzlich überliefert \textit{L} (Bl.\;\href{\Scan{LH_37_05_018r}}{18\,r\textsuperscript{o}}) einen gestrichenen, hier nicht wiedergegeben Entwurf zum Diagramm.%
}}%
}%
\end{edfigure}%
\pstart%
Hoc ut faciliore alio exemplo tentemus,%
\indexS{exemplum}
sit
\edtext{%
\edtext{%
libra \textit{DE}%
\indexS{libra}%
}{%
\lemma{libra \textit{DE}}%
\Cfootnote{%
Vgl. das Diagramm \bracks{\textit{Fig.\:2}}.%
}}
et brachium \textit{FG}%
\indexS{brachium librae}
instructum trochlea~\textit{H}%
\indexS{trochlea}
onerata ponderibus \textit{I},\;\textit{K},%
\indexS{pondus onerans trochleam}
sitque libra ita onerata %
\indexS{libra}
ut si ponderi
in lance \textit{E} posito%
\indexS{lanx librae}
addatur pondus \textit{I} %
pendens ex \textit{H},
fiat aequilibrium%
\indexS{aequilibrium librae}
cum pondere lancis \textit{D}.%
\indexS{lanx librae}%
}{%
\lemma{\textit{DE}}%
\Bfootnote{%
\textit{(1)}~in aequipondio ita constituta\indexS{aequipondium} %
\textit{(2)}~et brachium %
\bracks{...} trochlea \textit{H} %
\textit{(a)}~circa quam %
\textit{(b)}~onerata ponderibus %
\bracks{...} lancis \textit{D}.%
~\textit{L}%
}}
Hoc facto ajo%
\bracks{:}
si pondus \textit{K} praevaleat%
\indexS{pondus praevalens}
ponderi \textit{I},
\ledpb%
non nisi pondus \textit{I},
quippe minus%
\bracks{,}
considerari debere %
velut ex
\edtext{%
\bracks{trochlea}%
}{%
\lemma{brachio}%
\Bfootnote{%
\textit{L\;ändert Hrsg.}%
}}
\textit{H} ponderans,%
\indexS{brachium librae}
ac proinde manere aequilibrium%
\indexS{aequilibrium librae}%
\bracks{;}
si vero
\edtext{%
pondera \textit{K}\;et\;\textit{I}
in aequilibrio%
\indexS{aequilibrium ponderum}
essent,%
}{%
\lemma{pondera}%
\Bfootnote{%
\textit{(1)}~\textit{H}\;et\;\textit{K} %
\textit{(2)}~\textit{K}\;et\;\textit{I} %
\textit{(a)}~aequipa %
\textit{(b)}~in aequilibrio essent,%
~\textit{L}%
}}
utrumque pondus adjungendum esset ad pondus \textit{E},
atque ideo ipsum praevaleret%
\indexS{pondus praevalens}%
\bracks{;}
si vero pondus \textit{K} esset minus quam pondus \textit{I},
tunc pondus \textit{K} solum addendum esset ponderi \textit{E},
non vero
\edtext{%
pondus \textit{I},
adeoque%
}{%
\lemma{pondus}%
\Bfootnote{%
\textit{(1)}~\textit{D}, ade %
\textit{(2)}~\textit{I}, adeoque%
~\textit{L}%
}}
praevaleret lanx \textit{D}.%
\indexS{lanx librae}
Intelligitur
\edtext{%
autem brachium%
\indexS{brachium librae}%
}{%
\lemma{autem}%
\Bfootnote{%
\textit{(1)}~libram %
\textit{(2)}~brachium%
~\textit{L}%
}}
\textit{FG} non multum frictione%
\indexS{frictio chordae}
chordae \textit{KHI}%
\indexS{chorda}
affici.
Si vero frictio metueretur,%
\indexS{frictio chordae}
in promtu esset remedium elegans,%
\indexS{remedium frictioni}
scilicet si
\edtext{%
trochlea \textit{H}%
\indexS{trochlea}
non%
}{%
\lemma{trochlea}%
\Bfootnote{%
\bars{ista, \textit{gestr.}} \textit{H} non%
~\textit{L}%
}}
esset in brachio,%
\indexS{brachium librae}
sed in ligno%
\indexS{lignum natans in aqua}
quod nataret in aqua%
\indexS{aqua in vase}
qua vas brachio \textit{HG} sustentatum,
plenum est.%
\indexS{vas plenum aqua}
Qui modus videtur perelegans
efficiendi%
\indexS{modus efficiendi}
ne res subjecta sentiat frictionem.%
\indexS{frictio rei subjectae}
\pend%
\begin{edfigure}[1.750em][1.250em]%
\includegraphics%
{gesamttex/edit_VIII-4/images/LH_37_05_018_d2_018r.pdf}%
\edcaption[0.500em]{%
\bracks{\textit{Fig.\:2}}%
}%
\end{edfigure}%
\pstart%
Caeterum
\edtext{%
videtur hujus experimenti%
\indexS{experimentum}
successus%
\indexS{successus experimenti}%
}{%
\lemma{videtur}%
\Bfootnote{%
\textit{(1)}~hoc expe %
\textit{(2)}~hujus experimenti%
~\textit{L}%
}}
jam ex
\edtext{%
alio faciliore%
}{%
\lemma{alio}%
\Bfootnote{%
\textit{(1)}~facile %
\textit{(2)}~faciliore%
~\textit{L}%
}}
dijudicari posse.
\edlabel{LH_37_05_018_018r_Gdnkexp-1}%
\edtext{}{%
\xxref{LH_37_05_018_018r_Gdnkexp-1}{LH_37_05_018_018r_Gdnkexp-2}%
\lemma{Sit libra \bracks{...} sentiri}%
\Cfootnote{%
Auf das hier geschilderte Gedankenexperiment greift der spätere Textkomplex N.\;\ref{paradoxon_intro} \textit{Paradoxon de aeris gravitate} im vorliegenden Band zurück; siehe die editorische Einleitung dazu (S.\;\pageref{paradoxonbarometricum_intro}\,f.).%
}}%
Sit libra bilanx%
\indexS{libra bilanx}%
\indexS{bilanx}
aequilibrata,%
\indexS{libra aequilibrata}
ex uno ejus
\edtext{%
brachi\bracks{orum}%
\indexS{brachium librae}%
}{%
\lemma{brachiis}%
\Bfootnote{%
\textit{L\;ändert Hrsg.}%
}}
suspensum
vas aqua plenum;%
\indexS{vas plenum aqua}
jam adjiciatur pondus novum,%
\indexS{pondus adjectum}
quod in aqua%
\indexS{aqua in vase}
submergi potest;%
\indexS{pondus submersum}
ajo
non ante %
ponderis hujus adjectionem%
\indexS{adjectio ponderis} %
vim suam excituram%
\indexS{vis excita a pondere adjecto}
atque aequilibrium mutaturam,%
\indexS{aequilibrium mutatum a pondere adjecto}
quam ubi fundum attinget,%
\indexS{fundus vasis}
atque ita medio tempore non se faciet sentiri.%
\edlabel{LH_37_05_018_018r_Gdnkexp-2}
Quod ita evenire debere %
jam aliunde concludo%
\bracks{,}
ex familiari atque indubitato %
experimento%
\indexS{experimentum}%
\bracks{:}
si scilicet secus fieret,
\edtext{%
quoniam aquae%
\indexS{aqua in vase}%
}{%
\lemma{quoniam}%
\Bfootnote{%
\bars{scilicet \textit{gestr.}} aquae%
~\textit{L}%
}}
in vase%
\indexS{vas plenum aqua}
incumbit,
\ledpb%
tunc eodem jure etiam lapis%
\indexS{lapis incumbens aeri}
qui soli aeri incumbit%
\bracks{,}
stantem inferius lancem%
\indexS{lanx librae}
premeret %
etiam antequam
lapsu
ad eam perveniet;
quod non est.
Ergo
\edtext{%
nec lancem%
\indexS{lanx librae}%
}{%
\lemma{nec}%
\Bfootnote{%
\textit{(1)}~aquam %
\textit{(2)}~lancem%
~\textit{L}%
}}
inclinabit
si aquae incumbat.
Ratio autem est,%
\indexS{ratio}
\edtext{%
quia dum descendit%
}{%
\lemma{quia}%
\Bfootnote{%
\textit{(1)}~dum %
\textit{(2)}~dum %
\textit{(a)}~asce %
\textit{(b)}~descendit%
~\textit{L}%
}}
lapis in aqua,%
\indexS{lapis descendens in aqua}%
\indexS{aqua in vase}
attollit aquam paris voluminis,%
\indexS{aqua paris voluminis}
quae minus ponderat quam ipse,
minus ergo nempe
quia id
\edtext{%
\bracks{quod}%
}{%
\lemma{quod}%
\Bfootnote{%
\textit{erg. Hrsg.}%
}}
non descendit,
sed sustentatur,
solum ponderabit;
ac proinde
idem et in libra sicca%
\indexS{libra sicca}
eveniet,
eodem modo enim
pondus\;\textit{K} majus descendit,
minus\;\textit{I} ascendit,
solum ergo\;\textit{I} a bilance%
\indexS{bilanx}
sustentatur.
\pend%
\selectlanguage{latin}%
\frenchspacing%
\normalsize%
